% ====================================================================== Task 2B
\clearpage
\section{Task 2B: Dynamic I-RRIP}
I-RRIP helps Trace2 and Trace3 but hurts Trace1. Dynamic I-RRIP (\code{replacement/dyn\_irrip.l2c\_repl}) chooses
between SRRIP and I-RRIP while the program runs.

\subsection{Design}
The policy uses \textbf{set dueling}. A few \textbf{leader sets always run SRRIP} and, in the final design, the same
number of leader sets always run I-RRIP. All other sets (the followers) use whichever policy is currently winning.
Leader sets are placed by complement-select: the 1024 sets are split into $N$ groups of $1024/N$ sets, and in group
$k$ the set at offset $k$ is an SRRIP leader and the set at offset $1024/N-1-k$ is an I-RRIP leader. This spreads
the leaders over the whole index range, so they do not line up with power-of-two strides of the program.

The leader sets count instruction and data (non-instruction) \textbf{demand} accesses and misses; writebacks are not
counted. We implemented two decision rules (\code{DYN\_MODE}):
\begin{itemize}
  \item \textbf{Mode 0: the formula from the hint.} In the SRRIP leaders,
        ratio = (instr.\ accesses + instr.\ misses) / (total accesses + total misses).
        Every \code{DYN\_EPOCH} = 1024 leader accesses, the followers use I-RRIP if ratio $> T$. The four counters are
        then halved so that old program phases fade out. We measure the ratio in the SRRIP sets on purpose: in the
        I-RRIP sets, the policy we are judging would already have removed instruction misses and biased the measurement.
  \item \textbf{Mode 1: miss duel (submitted).} A 10-bit saturating counter PSEL starts at 512. A demand miss in an
        SRRIP leader adds its weight to PSEL, and a miss in an I-RRIP leader subtracts it. A data miss weighs 1 and an
        instruction miss weighs $\boldsymbol{W = 2}$. Every 1024 SRRIP-leader accesses, the followers use I-RRIP if
        PSEL $\ge 512$, i.e.\ if the I-RRIP leaders produced fewer weighted misses.
\end{itemize}

\subsection{The formula from the hint: threshold sensitivity}
\begin{table}[H]\centering\small
\begin{tabular}{@{}lcccccc@{}}
\toprule
Threshold & \makecell{Speedup\\Trace1} & \makecell{Speedup\\Trace2} & \makecell{Speedup\\Trace3} & \makecell{Geomean\\speedup}
          & \makecell{Fraction I-RRIP\\(T1/T2/T3)} & \makecell{Mean ratio\\(T1/T2/T3)} \\
\midrule
$T = 5\%$  & 0.9962 & 1.0265 & 1.0048 & 1.0091 & 1.00 / 1.00 / 1.00 & 0.44 / 0.61 / 0.31 \\
$T = 10\%$ & 0.9962 & 1.0231 & 1.0047 & 1.0079 & 1.00 / 0.98 / 1.00 & 0.44 / 0.61 / 0.31 \\
$T = 15\%$ & 0.9962 & 1.0230 & 1.0045 & 1.0079 & 1.00 / 0.98 / 0.86 & 0.44 / 0.61 / 0.31 \\
$T = 20\%$ & 0.9960 & 1.0226 & 0.9994 & 1.0059 & 1.00 / 0.96 / 0.57 & 0.44 / 0.61 / 0.31 \\
$T = 25\%$ & 0.9960 & 1.0226 & 0.9985 & 1.0056 & 1.00 / 0.96 / 0.43 & 0.44 / 0.61 / 0.31 \\
$T = 30\%$ & 0.9958 & 1.0219 & 0.9986 & 1.0053 & 1.00 / 0.94 / 0.43 & 0.44 / 0.61 / 0.31 \\
$T = 40\%$ & 0.9964 & 1.0218 & 0.9987 & 1.0056 & 0.85 / 0.92 / 0.29 & 0.44 / 0.61 / 0.31 \\
$T = 50\%$ & 0.9992 & 1.0211 & 0.9989 & 1.0063 & 0.04 / 0.90 / 0.29 & 0.44 / 0.61 / 0.31 \\
$T = 60\%$ & 0.9985 & 1.0182 & 0.9990 & 1.0052 & 0.00 / 0.82 / 0.00 & 0.45 / 0.61 / 0.31 \\
$T = 70\%$ & 0.9985 & 1.0027 & 0.9990 & 1.0001 & 0.00 / 0.00 / 0.00 & 0.45 / 0.61 / 0.31 \\
\bottomrule
\end{tabular}
\caption{Mode 0 (formula from the hint) with 10 thresholds; 32 SRRIP leader sets, epoch 1024. ``Fraction I-RRIP'' is
the share of ROI decisions that chose I-RRIP, ``mean ratio'' the average measured ratio.}
\label{tab:t2b-thr}
\end{table}

\pgfplotstableread{
T  t1     t2     t3     gm     f1   f2   f3
5  0.9962 1.0265 1.0048 1.0091 1.00 1.00 1.00
10 0.9962 1.0231 1.0047 1.0079 1.00 0.98 1.00
15 0.9962 1.0230 1.0045 1.0079 1.00 0.98 0.86
20 0.9960 1.0226 0.9994 1.0059 1.00 0.96 0.57
25 0.9960 1.0226 0.9985 1.0056 1.00 0.96 0.43
30 0.9958 1.0219 0.9986 1.0053 1.00 0.94 0.43
40 0.9964 1.0218 0.9987 1.0056 0.85 0.92 0.29
50 0.9992 1.0211 0.9989 1.0063 0.04 0.90 0.29
60 0.9985 1.0182 0.9990 1.0052 0.00 0.82 0.00
70 0.9985 1.0027 0.9990 1.0001 0.00 0.00 0.00
}\tThr
\begin{figure}[H]
\centering
\begin{subfigure}{0.49\textwidth}\centering
  \lineplot{\tThr}{T}{t1,t2,t3,gm}{xlabel={Threshold $T$ (\%)}, ylabel={Speedup over LRU}, xtick={5,10,...,30,40,50,60,70},
    ymin=0.994, ymax=1.03, speedline,
    legend entries={Trace1,Trace2,Trace3,Geomean}, legend columns=-1, legend to name=leg:thr}
  \caption{Speedup}
\end{subfigure}\hfill
\begin{subfigure}{0.49\textwidth}\centering
  \lineplot{\tThr}{T}{f1,f2,f3}{xlabel={Threshold $T$ (\%)}, ylabel={Fraction of decisions = I-RRIP},
    xtick={5,10,...,30,40,50,60,70}, ymin=-0.05, ymax=1.05,
    yticklabel style={/pgf/number format/.cd, fixed, fixed zerofill, precision=1}}
  \caption{Share of decisions that chose I-RRIP}
\end{subfigure}\\[3pt]
\ref*{leg:thr}
\caption{Mode 0: sensitivity to the threshold $T$.}
\label{fig:t2b-thr}
\end{figure}

The measured ratio hardly changes within a trace: about 0.44 on Trace1, 0.61 on Trace2 and 0.31 on Trace3.
Low thresholds ($T \le 15\%$) almost always choose I-RRIP and match static I-RRIP; $T = 70\%$ never chooses it and
equals SRRIP. In between, the traces switch to SRRIP in the order of their ratios: Trace3 first, then Trace1, then Trace2.

That order is the problem. \textbf{Trace3 gains from I-RRIP} ($+0.50\%$) even though only about 31\% of its leader
traffic is instruction traffic, and \textbf{Trace1 loses} ($-0.38\%$) even though 44\% is. A threshold between the two
ratios ($T$ = 30--40\%) gives Trace3 SRRIP and Trace1 I-RRIP, which is exactly backwards. $T = 50\%$ gets Trace1 right
but still takes I-RRIP away from Trace3. The best threshold is the lowest one ($T = 5\%$, $+0.91\%$), so the formula
works best when it does nothing.

The reason is that the ratio tells us \emph{how much} of the L2 traffic is instruction traffic, not \emph{whether}
protecting it pays off. That depends on what the protection costs. On Trace1 the data has L2 reuse, and moving data
insertion to RRPV 3 adds 2.57 data MPKI while removing only 1.30 instruction MPKI. To decide correctly, the policy has
to measure the effect of I-RRIP itself, and that is what set dueling against I-RRIP leader sets does.

\subsection{Miss duel: weight, leader sets and epoch}
\begin{table}[H]\centering\small
\begin{tabular}{@{}lccccc@{}}
\toprule
Instr.\ miss weight & \makecell{Speedup\\Trace1} & \makecell{Speedup\\Trace2} & \makecell{Speedup\\Trace3}
                    & \makecell{Geomean\\speedup} & \makecell{Fraction I-RRIP\\(T1/T2/T3)} \\
\midrule
$W = 1$         & 0.9993 & 1.0268 & 0.9994 & 1.0084 & 0.00 / 1.00 / 0.00 \\
$W = 1.5$       & 0.9993 & 1.0267 & 1.0012 & 1.0090 & 0.00 / 1.00 / 0.86 \\
\chosen{$W = 2$}& 0.9993 & 1.0268 & 1.0048 & 1.0102 & 0.00 / 1.00 / 1.00 \\
$W = 3$         & 0.9991 & 1.0267 & 1.0048 & 1.0101 & 0.04 / 1.00 / 1.00 \\
$W = 4$         & 0.9965 & 1.0267 & 1.0048 & 1.0092 & 0.88 / 1.00 / 1.00 \\
$W = 8$         & 0.9962 & 1.0267 & 1.0048 & 1.0092 & 1.00 / 1.00 / 1.00 \\
\bottomrule
\end{tabular}
\caption{Mode 1 with six instruction-miss weights $W$ (a data miss weighs 1; $W = 1.5$ is implemented as steps of 3
and 2), 32 + 32 leader sets, epoch 1024. Highlighted: submitted.}
\label{tab:t2b-w}
\end{table}

\pgfplotstableread{
W   t1     t2     t3     gm     f1   f2   f3
1   0.9993 1.0268 0.9994 1.0084 0.00 1.00 0.00
1.5 0.9993 1.0267 1.0012 1.0090 0.00 1.00 0.86
2   0.9993 1.0268 1.0048 1.0102 0.00 1.00 1.00
3   0.9991 1.0267 1.0048 1.0101 0.04 1.00 1.00
4   0.9965 1.0267 1.0048 1.0092 0.88 1.00 1.00
8   0.9962 1.0267 1.0048 1.0092 1.00 1.00 1.00
}\tW
\begin{figure}[H]
\centering
\begin{subfigure}{0.49\textwidth}\centering
  \lineplot{\tW}{W}{t1,t2,t3,gm}{xlabel={Instruction-miss weight $W$}, ylabel={Speedup over LRU}, xtick={1,1.5,2,3,4,8},
    xticklabel style={/pgf/number format/.cd, fixed}, ymin=0.994, ymax=1.03, speedline,
    legend entries={Trace1,Trace2,Trace3,Geomean}, legend columns=-1, legend to name=leg:w}
  \caption{Speedup}
\end{subfigure}\hfill
\begin{subfigure}{0.49\textwidth}\centering
  \lineplot{\tW}{W}{f1,f2,f3}{xlabel={Instruction-miss weight $W$}, ylabel={Fraction of decisions = I-RRIP},
    xtick={1,1.5,2,3,4,8}, xticklabel style={/pgf/number format/.cd, fixed}, ymin=-0.05, ymax=1.05,
    yticklabel style={/pgf/number format/.cd, fixed, fixed zerofill, precision=1}}
  \caption{Share of decisions that chose I-RRIP}
\end{subfigure}\\[3pt]
\ref*{leg:w}
\caption{Mode 1: sensitivity to the instruction-miss weight $W$.}
\label{fig:t2b-w}
\end{figure}

\newcommand{\cell}[2]{\makecell{#1\\[-1pt]{\scriptsize\color{ink2}#2}}}
\begin{table}[H]\centering\small
\begin{tabular}{@{}lccc@{}}
\toprule
Weight & 16 + 16 leader sets & 32 + 32 leader sets & 64 + 64 leader sets \\
\midrule
$W = 1$   & \cell{1.0104}{0.00 / 1.00 / 1.00 \cmark} & \cell{1.0084}{0.00 / 1.00 / 0.00} & \cell{1.0095}{1.00 / 1.00 / 1.00} \\
$W = 1.5$ & \cell{1.0104}{0.00 / 0.96 / 1.00 \cmark} & \cell{1.0090}{0.00 / 1.00 / 0.86 \cmark} & \cell{1.0095}{1.00 / 1.00 / 1.00} \\
$W = 2$   & \cell{1.0093}{0.93 / 1.00 / 1.00} & \cellcolor{hl}\cell{\textbf{1.0102}}{0.00 / 1.00 / 1.00 \cmark} & \cell{1.0095}{1.00 / 1.00 / 1.00} \\
\bottomrule
\end{tabular}
\caption{Geometric-mean speedup for every combination of weight and number of leader sets, with the share of decisions
that chose I-RRIP on Trace1 / Trace2 / Trace3 (\cmark: SRRIP on Trace1 and I-RRIP on Trace2 and Trace3, the best
static choice per trace). Highlighted: submitted.}
\label{tab:t2b-grid}
\end{table}

\begin{table}[H]\centering\small
\begin{tabular}{@{}lccccc@{}}
\toprule
\makecell[l]{Variant (mode 1, $W = 2$,\\32 + 32 leaders unless stated)} & \makecell{Speedup\\Trace1} & \makecell{Speedup\\Trace2}
  & \makecell{Speedup\\Trace3} & \makecell{Geomean\\speedup} & \makecell{Fraction I-RRIP\\(T1/T2/T3)} \\
\midrule
epoch 256                       & 0.9986 & 1.0268 & 1.0048 & 1.0100 & 0.00 / 1.00 / 1.00 \\
\chosen{epoch 1024 (submitted)} & 0.9993 & 1.0268 & 1.0048 & 1.0102 & 0.00 / 1.00 / 1.00 \\
epoch 4096                      & 0.9993 & 1.0267 & 1.0048 & 1.0102 & 0.00 / 1.00 / 1.00 \\
\midrule
$W = 1$, epoch 256              & 0.9986 & 1.0268 & 0.9994 & 1.0082 & 0.00 / 1.00 / 0.00 \\
$W = 1$, epoch 1024             & 0.9993 & 1.0268 & 0.9994 & 1.0084 & 0.00 / 1.00 / 0.00 \\
$W = 1$, epoch 4096             & 0.9993 & 1.0268 & 0.9994 & 1.0084 & 0.00 / 1.00 / 0.00 \\
\midrule
mode 0, $T = 25\%$, epoch 256   & 0.9957 & 1.0229 & 0.9986 & 1.0057 & 1.00 / 0.95 / 0.36 \\
mode 0, $T = 25\%$, epoch 1024  & 0.9960 & 1.0226 & 0.9985 & 1.0056 & 1.00 / 0.96 / 0.43 \\
mode 0, $T = 25\%$, epoch 4096  & 0.9962 & 1.0265 & 1.0048 & 1.0091 & 1.00 / 1.00 / 1.00 \\
\bottomrule
\end{tabular}
\caption{Epoch length: the followers re-read the decision every \emph{epoch} SRRIP-leader accesses.}
\label{tab:t2b-epoch}
\end{table}

\textbf{Weight.} With $W = 1$ (a plain miss count) and 32 + 32 leaders, the duel correctly picks SRRIP for Trace1 and
I-RRIP for Trace2, but it also picks SRRIP for Trace3, where the two policies differ by only a few misses. $W = 2$
picks I-RRIP for Trace3 as well and gives $+1.02\%$. $W = 3$ is a tie ($+1.01\%$), while $W = 4$ and 8 give
instruction misses too much weight and send Trace1 back to I-RRIP. $W = 2$ also has a physical meaning: an
instruction miss stalls fetch, while the core can overlap part of a data miss with other work, so an instruction miss
costs roughly twice as much.

\textbf{Leader sets and weight together (Table~\ref{tab:t2b-grid}).} The leaders are only a sample of the cache, and
on Trace1 and Trace3 the two policies are close. On Trace1, I-RRIP adds about 2.5 data MPKI and removes about 1.2
instruction MPKI, so with $W = 2$ the weighted misses of the two leader groups are almost equal. The winner then depends
on which sets happen to be sampled, which is why neighbouring cells of the grid disagree on Trace1 or Trace3. Three
settings choose the best policy on every trace (\cmark): 16 + 16 leaders with $W = 1$ or 1.5, and 32 + 32 leaders with
$W = 2$. They are within 0.02\% of each other, a tie. We submit \textbf{32 + 32 leaders with $\boldsymbol{W = 2}$}:
32 leader sets per policy (6.25\% of the 1024 sets together) is the size DIP uses for caches of this size, and it
halves the sampling noise of 16 sets (the standard error of a sampled miss rate falls with the square root of the
number of sets). 64 + 64 leaders run 12.5\% of the cache with a fixed policy, and in our runs they also chose I-RRIP
for Trace1. A wrong choice is cheap on these traces anyway: the worst cell of the grid is $+0.84\%$, still better
than SRRIP ($+0.02\%$) and close to static I-RRIP ($+0.95\%$), because the duel is only unsure where the two policies
are close.

\textbf{Epoch (Table~\ref{tab:t2b-epoch}).} PSEL is updated all the time; the epoch only sets how often the followers
read it. Trace3 sends few requests to the L2, so with an epoch of 1024 it makes only 7 decisions in 25M instructions.
That is still enough, because its preference never changes. For the duel, all three epochs give the same decisions
(within 0.02\%). For mode 0 the epoch changes which phases are sampled, but no epoch makes the formula choose the
right policies.

\subsection{Results}
\begin{table}[H]\centering\small
\begin{tabular}{@{}lccccc@{}}
\toprule
Configuration & IPC & Speedup & L2C load MPKI & L2 instr.\ load MPKI & L2 data load MPKI \\
\midrule
\grp{6}{Trace1}
LRU (baseline)      & 0.2321 & 1.0000 & 17.70 & 10.37 & 7.33 \\
SRRIP               & 0.2320 & 0.9996 & 17.61 & 10.55 & 7.06 \\
I-RRIP              & 0.2312 & 0.9962 & 18.88 & 9.25  & 9.63 \\
\chosen{Dyn-I-RRIP} & 0.2320 & 0.9993 & 17.68 & 10.54 & 7.14 \\
\midrule
\grp{6}{Trace2}
LRU (baseline)      & 0.2499 & 1.0000 & 28.64 & 20.14 & 8.50 \\
SRRIP               & 0.2504 & 1.0020 & 27.27 & 19.32 & 7.95 \\
I-RRIP              & 0.2568 & 1.0276 & 23.25 & 12.33 & 10.92 \\
\chosen{Dyn-I-RRIP} & 0.2566 & 1.0268 & 23.33 & 12.53 & 10.79 \\
\midrule
\grp{6}{Trace3}
LRU (baseline)      & 0.3820 & 1.0000 & 3.96 & 1.12 & 2.84 \\
SRRIP               & 0.3816 & 0.9990 & 4.04 & 1.20 & 2.85 \\
I-RRIP              & 0.3839 & 1.0050 & 3.90 & 0.73 & 3.17 \\
\chosen{Dyn-I-RRIP} & 0.3838 & 1.0048 & 3.91 & 0.75 & 3.16 \\
\midrule
\grp{6}{Average (geometric mean for IPC and speedup, arithmetic mean otherwise)}
LRU (baseline)      & 0.2809 & 1.0000 & 16.77 & 10.54 & 6.22 \\
SRRIP               & 0.2809 & 1.0002 & 16.31 & 10.36 & 5.95 \\
I-RRIP              & 0.2836 & 1.0095 & 15.34 & 7.44  & 7.91 \\
\chosen{Dyn-I-RRIP} & 0.2838 & 1.0102 & 14.97 & 7.94  & 7.03 \\
\bottomrule
\end{tabular}
\caption{Task 2B: IPC, speedup over LRU and L2 load MPKIs (no prefetcher). Dyn-I-RRIP = mode 1, $W = 2$,
32 + 32 leaders, epoch 1024.}
\label{tab:t2b-res}
\end{table}

\replfig{lru,srrip,irrip,dyn}{LRU (baseline),SRRIP,I-RRIP,Dyn-I-RRIP}{leg:t2b}{4pt}
  {Task 2B: speedup and L2C / L2 instruction / L2 data load MPKI.}{fig:t2b-res}

Dynamic I-RRIP keeps the gains of I-RRIP where they exist (Trace2 $+2.68\%$, Trace3 $+0.48\%$) and falls back to SRRIP
on Trace1 ($-0.07\%$ instead of $-0.38\%$ with static I-RRIP). As a result it beats both SRRIP and I-RRIP in geometric
mean ($+1.02\%$ vs $+0.02\%$ and $+0.95\%$), and Trace1's data load MPKI goes back to the SRRIP level (7.14).

\textbf{Storage.} On top of SRRIP (2-bit RRPV per line, 16384 bits), the policy needs the instruction bit per line
(8192 bits), a 10-bit PSEL, a 10-bit epoch counter and the 1-bit follower mode: 8213 bits $\approx$ 1.0\,KiB in total.
Leader sets are recognised by comparing two 5-bit fields of the set index, so no table is needed. (Mode 0 would need
four 16-bit counters, 64 bits, instead of PSEL.)
